%%% Tento soubor obsahuje definice různých užitečných maker a prostředí %%%
%%% Další makra připisujte sem, ať nepřekáží v ostatních souborech.     %%%
%%% This file contains definitions of various useful macros and environments      %%%
%%% Assign additional macros here so that they do not interfere with other files. %%%

% pdfx načte xcolor se zastaralou volbou hyperref --- varování je neškodné
\usepackage{silence}
\WarningFilter{xcolor}{Package option `hyperref'}

\usepackage[a-2u]{pdfx}     % výsledné PDF bude ve standardu PDF/A-2u
                            % resulting PDF will be in the PDF / A-2u standard

\usepackage{ifpdf}
\usepackage{ifxetex}
\usepackage{ifluatex}

%% Přepneme na českou sazbu, fonty Latin Modern a kódování češtiny
\ifthenelse{\boolean{xetex}\OR\boolean{luatex}}
   { % use fontspec and OpenType fonts with utf8 engines
			\usepackage[english,czech]{babel}
			\usepackage[autostyle,english=british,czech=quotes]{csquotes}
			\usepackage{fontspec}
			\defaultfontfeatures{Ligatures=TeX,Scale=MatchLowercase}
   }
   {
			\usepackage[english,czech]{babel}
			\usepackage{lmodern}
			\usepackage[T1]{fontenc}
			\usepackage{textcomp}
			\usepackage[utf8]{inputenc}
			\usepackage[autostyle,english=british,czech=quotes]{csquotes}
	 }
\ifluatex
\makeatletter
\let\pdfstrcmp\pdf@strcmp
\makeatother
\fi

%%% Nastavení pro použití samostatné bibliografické databáze (až po inputenc/csquotes u pdf\TeX).
%%% Settings for using a separate bibliographic database.
\input{biblatex-setup}

%%% Další užitečné balíčky (jsou součástí běžných distribucí LaTeXu)
\usepackage{amsmath}        % rozšíření pro sazbu matematiky / extension for math typesetting
\usepackage{amsfonts}       % matematické fonty / mathematical fonts
\usepackage{amssymb}        % symboly / symbols
\usepackage{amsthm}         % sazba vět, definic apod. / typesetting of sentences, definitions, etc.
\usepackage{bm}             % tučné symboly (příkaz \bm) / bold symbols (\bm command)
\usepackage{graphicx}       % vkládání obrázků / graphics inserting
% nopatch=item: enumitem mění \@item — microtype by jinak hlásilo „Unable to apply patch item“
\usepackage[nopatch=item]{microtype} % jemnější zalamování a kerning (pdflatex)
\usepackage{listings}       % vylepšené prostředí pro strojové písmo / improved environment for source codes typesetting
\usepackage{fancyhdr}       % prostředí pohodlnější nastavení hlavy a paty stránek / environment for more comfortable adjustment of the head and foot of the pages
\usepackage{icomma}         % inteligetní čárka v matematickém módu / intelligent comma in math mode
\usepackage{dcolumn}        % lepší zarovnání sloupců v tabulkách / better alignment of columns in tables
\usepackage{booktabs}       % lepší vodorovné linky v tabulkách / better horizontal lines in tables
\usepackage{tabularx}       % tabulky s automatickou šířkou sloupců / tables with auto-width columns
\usepackage{float}          % přesné umístění floatů ([H]) / exact float placement
\makeatletter
\@ifpackageloaded{xcolor}{
   \@ifpackagewith{xcolor}{usenames}{}{\PassOptionsToPackage{usenames}{xcolor}}
  }{\usepackage[usenames,table]{xcolor}} % barevná sazba / color typesetting
\makeatother
\usepackage{colortbl}       % \cellcolor v tabulkách / cell background colors
\usepackage{tikz}           % diagramy / diagrams
\usetikzlibrary{positioning,calc}
\usepackage{pgfplots}       % grafy s daty / data plots
\pgfplotsset{compat=1.18}
\usepackage{multicol}       % práce s více sloupci na stránce / work with multiple columns on a page
\usepackage{pdflscape}      % otočení stránky na šířku (landscape) / landscape pages
\usepackage{caption}
\usepackage{enumitem}
\usepackage{needspace}
\setlist[itemize]{noitemsep, topsep=0pt, partopsep=0pt}
\setlist[enumerate]{noitemsep, topsep=0pt, partopsep=0pt}
\setlist[description]{noitemsep, topsep=0pt, partopsep=0pt}

\usepackage{tocloft}
\setlength\cftparskip{0pt}
\setlength\cftbeforechapskip{1.5ex}
\setlength\cftfigindent{0pt}
\setlength\cfttabindent{0pt}
\setlength\cftbeforeloftitleskip{0pt}
\setlength\cftbeforelottitleskip{0pt}
\setlength\cftbeforetoctitleskip{0pt}
\renewcommand{\cftlottitlefont}{\Huge\bfseries\sffamily}
\renewcommand{\cftloftitlefont}{\Huge\bfseries\sffamily}
\renewcommand{\cfttoctitlefont}{\Huge\bfseries\sffamily}

% vyznaceni odstavcu
% differentiation of new paragraphs
\parindent=0pt
\parskip=11pt

% zakaz vdov a sirotku - jednoradkovych pocatku ci koncu odstavcu na prechodu mezi strankami
% Prohibition of widows and orphans - single-line beginning and end of paragraph at the transition between pages
\clubpenalty=1000
\widowpenalty=1000
\displaywidowpenalty=1000

% nastaveni radkovani
% setting of line spacing
\renewcommand{\baselinestretch}{1.20}

% nastaveni pro nadpisy - tucne a bezpatkove (titlesec místo zastaralého sectsty)
% settings for headings - bold and sans serif
\usepackage{titlesec}
\titleformat{\section}{\sffamily\Large\bfseries}{\thesection}{1em}{}
\titleformat{\subsection}{\sffamily\large\bfseries}{\thesubsection}{1em}{}
\titleformat{\subsubsection}{\sffamily\normalsize\bfseries}{\thesubsubsection}{1em}{}
\titleformat{\paragraph}{\sffamily\normalsize\bfseries}{\theparagraph}{1em}{}
\titleformat{\subparagraph}{\sffamily\normalsize\bfseries}{\thesubparagraph}{1em}{}

% nastavení hlavy a paty stránek
% page head and foot settings
\makeatletter
\if@twoside%
    \fancypagestyle{fancyx}{%
			\fancyhf{}                                     
      \fancyhead[RE]{\rightmark}                  
      \fancyhead[LO]{\leftmark}                  
      \fancyfoot[RO,LE]{\thepage}                    
      \renewcommand{\headrulewidth}{.5pt}            
      \renewcommand{\footrulewidth}{.5pt}            
    }
    \fancypagestyle{plain}{%
			\fancyhf{}                                     
    	\fancyfoot[RO,LE]{\thepage}                    
    	\renewcommand{\headrulewidth}{0pt}             
    	\renewcommand{\footrulewidth}{0.5pt}
    }          
\else
    \fancypagestyle{fancyx}{%
			\fancyhf{}                                     
      \fancyhead[R]{\leftmark}                  
      \fancyfoot[R]{\thepage}                    
      \renewcommand{\headrulewidth}{.5pt}            
      \renewcommand{\footrulewidth}{.5pt}            
    }
    \fancypagestyle{plain}{%                       
    	\fancyhf{} % clear all header and footer fields
    	\fancyfoot[R]{\thepage}                    
    	\renewcommand{\headrulewidth}{0pt}             
    	\renewcommand{\footrulewidth}{0.5pt}
    }          
\fi
\renewcommand*{\cleardoublepage}{\clearpage\if@twoside \ifodd\c@page\else
	\hbox{}%
	\thispagestyle{empty}%
	\newpage%
	\if@twocolumn\hbox{}\newpage\fi\fi\fi
}
\makeatother

% Tato makra přesvědčují mírně ošklivým trikem LaTeX, aby hlavičky kapitol
% sázel příčetněji a nevynechával nad nimi spoustu místa. Směle ignorujte.
% These macros convince with a slightly ugly LaTeX trick to make chapter headers
% bet more sane and didn't miss a lot of space above them. Be boldly ignore it.
\makeatletter
\def\@makechapterhead#1{
  {\parindent \z@ \raggedright \sffamily
   \Huge\bfseries \thechapter. #1
   \par\nobreak
   \vskip 20\p@
}}
\def\@makeschapterhead#1{
  {\parindent \z@ \raggedright \sffamily
   \Huge\bfseries #1
   \par\nobreak
   \vskip 20\p@
}}
\makeatother

% Trochu volnější nastavení dělení slov, než je default.
% Slightly looser hyphenation setting than default.
\lefthyphenmin=2
\righthyphenmin=2

% Povolí LaTeXu roztáhnout mezery při nouzi, místo přetečení řádku.
% Standardní řešení pro jazyky s~dlouhými slovy (čeština, němčina).
\emergencystretch=3em

% Zapne černé "slimáky" na koncích řádků, které přetekly, abychom si jich lépe všimli.
% Turns on the black "snails" at the ends of the lines that overflowed to get us noticed them better.
\overfullrule=1mm

%% Balíček hyperref, kterým jdou vyrábět klikací odkazy v PDF,
%% ale hlavně ho používáme k uložení metadat do PDF (včetně obsahu).
%% Většinu nastavítek přednastaví balíček pdfx.
%% A hyperref package that can be used to produce clickable links in PDF,
%% but we mainly use it to store metadata in PDF (including content).
%% Most settings are preset by the pdfx package.
\hypersetup{unicode}
\hypersetup{breaklinks=true}
\hypersetup{hidelinks}

\renewcommand{\UrlBreaks}{\do\/\do\=\do\+\do\-\do\_\do\ \do\a\do\b\do\c\do\d%
\do\e\do\f\do\g\do\h\do\i\do\j\do\k\do\l\do\m\do\n\do\o\do\p\do\q\do\r\do\s%
\do\t\do\u\do\v\do\w\do\x\do\y\do\z\do\A\do\B\do\C\do\D\do\E\do\F\do\G\do\H%
\do\I\do\J\do\K\do\L\do\M\do\N\do\O\do\P\do\Q\do\R\do\S\do\T\do\U\do\V\do\W%
\do\X\do\Y\do\Z\do\1\do\2\do\3\do\4\do\5\do\6\do\7\do\8\do\9\do\0}

%%% Prostředí pro sazbu kódu, případně vstupu/výstupu počítačových
%%% programů. (Vyžaduje balíček listings -- fancy verbatim.)
%%% Environment for source code typesetting, or computer input/output
%%% programs. (Requires package listings - fancy verbatim.)
\lstnewenvironment{code}{\lstset{basicstyle=\small, frame=single}}{}

%%% Prostředí pro označení stavu obsahu (workflow: raw → draft → finální)
%%% Environments for marking content status (workflow: raw → draft → final)
%%% Konvence: [RAW]/[DRAFT] je vždy na vlastním řádku
%%% Flagy \hideraw / \hidedraft se nastavují z příkazové řádky (viz Makefile)
\usepackage{comment}
\definecolor{rawcolor}{RGB}{120, 80, 120}      % tlumená fialová
\definecolor{draftcolor}{RGB}{80, 100, 120}    % tlumená modro-šedá

%%% Systém odkazování metrik (acro + custom příkazy)
%%% Dvě dimenze: CO měříme (P/Q/E) a JAK měříme (deter./kvalit./záznam.)
\usepackage{acro}
\usepackage{etoolbox}

\acsetup{
  first-style = short-long,   % P1 (issues before code)
  make-links  = true           % klikatelné → seznam zkratek
}
% Workaround: acro 3.8 ignoruje subsequent-style=short-long.
% Předefinujeme \ac aby vždy ukazoval plný tvar (= \acf).
\ExplSyntaxOn
\RenewDocumentCommand \ac {sO{}m} {
  \IfBooleanTF {#1}
    { \acf*[#2]{#3} }
    { \acf[#2]{#3} }
}
\ExplSyntaxOff
% \ac{P1}  → vždy "P1 (issues before code)" s hyperlinkem
% \acs{P1} → jen "P1" s hyperlinkem (pro opakování v těsné blízkosti)

% --- Procesní metriky (P) ---
\DeclareAcronym{P1}{short=P1, long=issues before code, tag=process}
\DeclareAcronym{P2}{short=P2, long=branch per issue, tag=process}
\DeclareAcronym{P3}{short=P3, long=test-first commits, tag=process}
\DeclareAcronym{P4}{short=P4, long=PRs linked to issues, tag=process}
\DeclareAcronym{P5}{short=P5, long=no existing test modifications, tag=process}
\DeclareAcronym{P6}{short=P6, long=commit message quality, tag=process}
\DeclareAcronym{P7}{short=P7, long=issue description quality, tag=process}
\DeclareAcronym{P8}{short=P8, long=PR description quality, tag=process}

% --- Produktové metriky (Q) ---
\DeclareAcronym{Q1}{short=Q1, long=API contract match, tag=quality}
\DeclareAcronym{Q2}{short=Q2, long=referenční test pass rate, tag=quality}
\DeclareAcronym{Q3}{short=Q3, long=mutation score, tag=quality}
\DeclareAcronym{Q4}{short=Q4, long=AC coverage, tag=quality}
\DeclareAcronym{Q5}{short=Q5, long=lint warnings, tag=quality}
\DeclareAcronym{Q6}{short=Q6, long=typecheck errors, tag=quality}
\DeclareAcronym{Q7}{short=Q7, long=složitost kódu, tag=quality}
\DeclareAcronym{Q8}{short=Q8, long=design quality, tag=quality}

% --- Metriky efektivity (E) ---
\DeclareAcronym{E1}{short=E1, long=vstup/výstup/cache tokeny, tag=efficiency}
\DeclareAcronym{E2}{short=E2, long=trvání, tag=efficiency}
\DeclareAcronym{E3}{short=E3, long=stabilita session, tag=efficiency}

% --- Skupiny: CO měříme ---
% Použití: \mgrp{P}, \mgrp{Q}, \mgrp{E}
\newcommand{\mgrp}[1]{%
  \ifstrequal{#1}{P}{\hyperref[tab:metriky-prehled]{procesní metriky (P1--P8)}}{}%
  \ifstrequal{#1}{Q}{\hyperref[tab:metriky-prehled]{produktové metriky (Q1--Q8)}}{}%
  \ifstrequal{#1}{E}{\hyperref[tab:metriky-prehled]{metriky efektivity (E1--E3)}}{}%
}

% --- Skupiny: JAK měříme ---
% Použití: \mmet{det}, \mmet{qual}, \mmet{zaz}
\newcommand{\mmet}[1]{%
  \ifstrequal{#1}{det}{\hyperref[tab:metriky-prehled]{deterministické metriky}}{}%
  \ifstrequal{#1}{qual}{\hyperref[tab:metriky-prehled]{kvalitativní metriky}}{}%
  \ifstrequal{#1}{zaz}{\hyperref[tab:metriky-prehled]{záznamové metriky}}{}%
}

\ifdefined\hideraw
  \excludecomment{raw}
\else
  \newenvironment{raw}{\par\noindent\textcolor{rawcolor}{[RAW]}\par\color{rawcolor}}{\par}
\fi

\ifdefined\hidedraft
  \excludecomment{draft}
\else
  \newenvironment{draft}{\par\noindent\textcolor{draftcolor}{[DRAFT]}\par\color{draftcolor}}{\par}
\fi
